\chapter*{Abstract}

\textbf{CodeCraft} is a domain-specific language (DSL) designed to enable children aged 10--15 
to learn programming by scripting actions within the Minecraft game environment. 
The project was developed by Chiril Boboc, Vasile Brînză, Cristian Bruma, Gabriela Bîtca, 
and Teodor Strulea, students of Software Engineering at the Technical University of Moldova.

The motivation for this project stems from the disconnect between traditional programming 
education and the interests of young learners. Existing tools either rely on abstract block-based 
environments or demand technical knowledge beyond the reach of beginners. CodeCraft addresses 
this gap by providing a text-based language with minimal syntax, immediate visual feedback, 
and a meaningful context rooted in a game children already engage with.

The report covers a domain analysis examining the educational context, target user research, 
and the analysis of the solution, followed by the language design chapter, which details the DSL requirements, 
formal grammar, and the implementation of the lexer and parser components.

\textbf{Keywords: } domain-specific language, programming education, gamification, Minecraft, 
computational thinking.
